\chapter{The State Machine}
\label{ch:states}

\section{A Mealy machine over branch pairs}

Two asks put side by side form a branch pair. Each cycle the pair is read as a state, and each
change of state from cycle \(N-1\) to cycle \(N\) is given a label. A machine whose output is
written on its transitions, a function of the state left and the input that moved it, is a Mealy
machine \cite{src:Mealy}. The language's clock is that output: a stream of states and labels from
which the health of a system can be read, as from a status page.

\section{The states}

Four states come from the two bits alone, and four from what the bits were measured with.

\begin{table}[h]
\centering
\small
\begin{tabular}{@{}l l p{6.4cm}@{}}
\toprule
state & reading & field term \\
\midrule
DUAL & 1, 1 & both sides healthy; active-active \\
LEAD & 1, 0 & left serves; active-passive \\
RITE & 0, 1 & right serves; active-passive \\
VOID & 0, 0 & outage \\
BUSY & both held, both dear & saturation \\
WAIT & one held, one past the timeout & latency \\
BLOK & a side ended its asker & error \\
GRAY & a side not asked & unknown; Belnap's None \\
\bottomrule
\end{tabular}
\caption{The base states.}
\label{tab:states}
\end{table}

BUSY, WAIT and BLOK are three of the four golden signals of production monitoring: saturation,
latency and errors \cite{src:SRE}. The fourth, traffic, is the rate of asks, which the protocol
sets and does not read.

GRAY is no answer, and it is not VOID, which is an answer of no. Belnap's four-valued logic for
a machine that answers questions has a value for exactly this, None, told nothing, beside True,
False and Both \cite{src:Belnap}. In his order of information None sits below every other value:
it is compatible with each of them until an ask is put, the language's ``every state at
once''. Entering GRAY is the loss of the ability to ask, labeled FUZZ, and leaving it is a
first observation, labeled FIZZ.

\begin{definition}[Reading a state off asks]
A side reads 1 where its ask held inside the timeout. A pair reads BLOK where a side ended its
asker, GRAY where a side was not asked, WAIT where one side held and the other held past the
timeout, and BUSY where both held and each cost more than every held ask of the baseline. The
edge BUSY is read against is the baseline's largest held cost, never a figure written in.
\end{definition}

The BUSY edge is the rank rule of Chapter~\ref{ch:protocol} applied to one side: a cost past the
baseline's least by more than the baseline's spread is past its largest.

\section{The transitions}

Table~\ref{tab:matrix} is the language's first transition table, the label for each pair
of past and present state.

\begin{table}[h]
\centering
\footnotesize
\setlength{\tabcolsep}{3.5pt}
\begin{tabular}{@{}l *{8}{c}@{}}
\toprule
past & DUAL & LEAD & RITE & VOID & BUSY & WAIT & BLOK & GRAY \\
\midrule
DUAL & NEXUS & DROP & SYNC & DROP & SYNC & SYNC & HALT & FUZZ \\
LEAD & JOIN & CORE & PASS & DROP & SYNC & SYNC & HALT & FUZZ \\
RITE & JOIN & BACK & SURV & DROP & SYNC & SYNC & HALT & FUZZ \\
VOID & SPRK & WAKE & WAKE & ZERO & SYNC & SYNC & HALT & FUZZ \\
BUSY & SYNC & SYNC & SYNC & DROP & DRAG & SYNC & JAMM & FUZZ \\
WAIT & SYNC & SYNC & SYNC & LOSS & SYNC & HOLD & HALT & FUZZ \\
BLOK & SYNC & SYNC & SYNC & DROP & SYNC & SYNC & DEAD & FUZZ \\
GRAY & FIZZ & FIZZ & FIZZ & FIZZ & FIZZ & FIZZ & FIZZ & \textendash \\
\bottomrule
\end{tabular}
\caption{The state transition table: rows are the state at cycle \(N-1\), columns the state at
cycle \(N\).}
\label{tab:matrix}
\end{table}

The diagonal, NEXUS, CORE, SURV, ZERO, DRAG, HOLD and DEAD, is the steady states. PASS and BACK
are the handoff between the two sides, failover and failback. DROP and JAMM are collapses: a
stable state lost, and saturation turned into refusal, which a distributed system knows as
gridlock. HALT is any state reaching a refusal.

A second table splits the transitions out of high load (Table~\ref{tab:energy}).

\begin{table}[h]
\centering
\footnotesize
\begin{tabular}{@{}l c c c c@{}}
\toprule
past & DUAL & LEAD & RITE & VOID \\
\midrule
BUSY & SURG & VENT & VENT & DROP \\
WAIT & SNAP & LEAD & RITE & LOSS \\
DUAL & NEXUS & TILT & TILT & FUSE \\
\bottomrule
\end{tabular}
\caption{The high load transition table.}
\label{tab:energy}
\end{table}

SURG is saturation clearing into a burst of throughput. VENT is a saturated pair dropping one
side, load shedding. SNAP is a lagging side catching up. TILT is DUAL decaying to one side. FUSE
is both sides lost at once with no decay: the breaker trips. The language leaves FUSE, JAMM and
HALT without a recovery path; the field's answer is the circuit breaker, which opens on failure,
stops sending, and after a wait lets one trial request through before it closes again
\cite{src:Nygard}. SPRK, from VOID straight to DUAL, is that trial succeeding.

\section{Superstates}

SYNC covers 22 cells of Table~\ref{tab:matrix}. A label that covers many transitions loses them:
a branch that reads only SYNC cannot tell DUAL to BUSY from BLOK to LEAD. The language's answer
is to make SYNC a superstate. A transition the table marks SYNC is written as a passage through
it, entering by FUZZ carrying the state it left and leaving by FIZZ carrying the state it reaches:
\[
  [\mathrm{DUAL}] \xrightarrow{\ \mathrm{FUZZ+DUAL}\ } [\mathrm{SYNC}]
  \xrightarrow{\ \mathrm{FIZZ+BUSY}\ } [\mathrm{BUSY}].
\]
GRAY is entered and left the same way. This is Harel's superstate \cite{src:Harel}: one state that
stands for a set of others, entered and left through labeled transitions. Here the labels carry
the identity of the transition, and the pair is read back from them exactly.

\begin{proposition}
\label{prop:atomic}
A trip through GRAY is atomic: a FUZZ opens it, one FIZZ closes it, and no FUZZ opens inside
another.
\end{proposition}

\begin{proof}
FIZZ labels only transitions out of GRAY to another state, FUZZ only transitions into GRAY from
another state, and GRAY to GRAY carries no label. Between a FUZZ and the next FIZZ the pair is
GRAY throughout, and no transition out of GRAY is a FUZZ.
\end{proof}

Over 500 random traces of 40 cycles, every side drawn from held, heavy, not held, late, ended and
unasked, every GRAY trip is atomic, and the exact pair of all 3,035 SYNC passages is read back
from their labels alone.

\section{A label is decided in its situation}

Where both tables name one transition they can disagree. DUAL to LEAD is DROP in the first table
and TILT in the second; DUAL to VOID is DROP in the first and FUSE in the second. The language
does not pick one for every case. Two names for one pair are two candidates, and the asks that
read the transition choose between them the way the gate chooses an arrangement: a bit excludes
what does not hold, and a magnitude ranks the survivors.

The protocol already carries what separates them. A side reading 0 carries why: past its
timeout, not held, or ended. DUAL to LEAD is TILT where the right side came in past its timeout
by a fraction of a cycle, and DROP where it did not answer at all or came in late by a timeout's
measure. DUAL to VOID is FUSE where both sides go at once with no decay, and DROP where the
collapse ran through a timeout. Where the line between a fraction and a timeout falls is read off
the baseline and never written in. Under this rule every pair the tables name takes a label from
its own candidates in each of the 256 situations two sides can be in.

\section{The tables read against asks}

Read against real asks, the tables raise six points of definition, each left for the language's
design to settle.
\begin{itemize}
\item TILT and WAIT name one event. Read off asks, a held side beside one a fraction late is
  WAIT, and DUAL to LEAD by TILT never arrives.
\item A name stands for a state and for a transition. The second table labels WAIT to LEAD as
  LEAD; CORE is the pair 1, 0 in the branch pair table and LEAD to LEAD in the transition table.
\item VOID does not separate two sides both past the timeout from two that both did not hold.
  The kind each ask carries does, and the record keeps it.
\item Two prose rules carry exceptions the tables hold: any state to BLOK is HALT, except BUSY to
  BLOK (JAMM) and BLOK to BLOK (DEAD); any stable state to VOID is DROP, except DUAL to VOID where
  both sides go at once (FUSE). The tables are taken as the definition.
\item DUAL at maximum is called unstable, and DUAL to DUAL, NEXUS, is listed as a steady state.
\item SHIFT, in the branch pair table, has five letters where every other keyword has four.
\end{itemize}
